\clearpage
\begin{summarybox}
\textbf{\textcolor{CSummary}{本结论真正值得掌握的是一个“垂直关系”，而不只是一个分式。}}
\[
\boxed{\frac{x_0}{a^2}\Delta x+\frac{y_0}{b^2}\Delta y=0
\quad\Longrightarrow\quad
\frac{x_0}{a^2}(x-x_0)+\frac{y_0}{b^2}(y-y_0)=0\ (M\ne O).}
\]
对 $y_0\ne0$ 的非退化弦，\textbf{最快速写}：
\[
\boxed{\qquad k=-\frac{b^2x_0}{a^2y_0}.\qquad}
\]
\textbf{\textcolor{CSummary}{三种核心证明}}
\begin{itemize}[leftmargin=*,itemsep=0.25em]
\item \textbf{最快：}两端点代入\,$\to$\,两式相减\,$\to$\,平方差\,$\to$\,中点之和\,$\to$\,斜率；适合考试速写。
\item \textbf{最低门槛：}中点对称，$A=(x_0+t,y_0+kt),\ B=(x_0-t,y_0-kt)$，相减消去 $t^2$；适合理解“从哪想到”。
\item \textbf{常规：}逐项写出两端点方程、相减、代入中点、分情形讨论；适合完整证明。
\end{itemize}
\textbf{补充直觉：}椭圆拉伸成圆，圆的弦中点半径垂直于弦，故出现加权垂直条件。
\medskip
\textbf{\textcolor{CSummary}{三种边界}}
\begin{itemize}[leftmargin=*,itemsep=0.25em]
\item $x_0=0,y_0\ne0$：\textbf{水平弦} $y=y_0$。
\item $y_0=0,x_0\ne0$：\textbf{竖直弦} $x=x_0$，斜率不存在。
\item $M=O$：\textbf{中心弦不唯一}，不可用斜率公式确定方向。
\end{itemize}
\textbf{使用前：}检查 $x_0^2/a^2+y_0^2/b^2<1$；这是\textbf{真实非退化弦的存在性条件}。
\medskip
\textbf{\textcolor{CSummary}{迁移：}}统一写 $x^2/A+y^2/B=1$：
$\frac{x_0}{A}\Delta x+\frac{y_0}{B}\Delta y=0$。无须死记“焦点在 $x$ 轴 / $y$ 轴”两套斜率公式。
\end{summarybox}
